\begin{lemma}
Let \(K\) be a finite field of order \(q\), let \(t\ge 2\), and let \(G(t,K)\) be the polarity graph. Then \(G(t,K)\) is an \((n,d,\lambda)\)-graph with
\[
  n=\frac{q^{t+1}-1}{q-1},\qquad
  d=\frac{q^t-1}{q-1},\qquad
  \lambda=\sqrt{d-\frac{q^{t-1}-1}{q-1}}.
\]
\end{lemma}
\begin{proof}
SKETCH:
The projective space has one point for each nonzero vector modulo multiplication by a nonzero scalar, so the number of points is \((q^{t+1}-1)/(q-1)\). For a fixed nonzero vector \(x\), the orthogonal hyperplane \(x^\perp\) has \(q^t\) vectors, one of which is zero, so the degree is \((q^t-1)/(q-1)\). If \(x\) and \(y\) represent distinct projective points, then they are linearly independent, and the common orthogonal space \(x^\perp\cap y^\perp\) has \(q^{t-1}\) vectors, giving \(a=(q^{t-1}-1)/(q-1)\) common neighbours. Therefore the adjacency matrix \(A\) satisfies \(A^2=aJ+(d-a)I\). On the all-ones vector \(A\) has eigenvalue \(d\); on the orthogonal complement of the all-ones vector, \(J\) vanishes, so all remaining eigenvalues have absolute value \(\sqrt{d-a}\).
\end{proof}
